\documentclass[10pt,twocolumn]{article}
\usepackage[margin=2cm]{geometry}
\usepackage{amsmath,amssymb}
\usepackage{booktabs}
\usepackage{graphicx}
\usepackage{xcolor}
\usepackage{multirow}
\usepackage{array}
\usepackage[hidelinks]{hyperref}
\usepackage{natbib}
\usepackage{caption}
\captionsetup{font=small}

\newcommand{\ie}{\emph{i.e.}}
\newcommand{\eg}{\emph{e.g.}}
\newcommand{\cf}{\emph{cf.}}
\newcommand{\etal}{\emph{et al.}}

\title{Instruction Routing in Vision-Language-Action Models:\\A Cross-Architecture Token-Level Analysis}

\author{
  Anonymous Authors\\
  \textit{Under Review at CoRL 2026}
}

\date{}

\begin{document}
\maketitle

% ============================================================
\begin{abstract}
Vision-Language-Action models (VLAs) fine-tune large language models on robot demonstration data, yet how training reshapes instruction routing and information flow at the token level remains uncharacterized---and whether such effects generalize across architectures is unknown.
We present a systematic experimental framework for token-level analysis of decoder-only VLAs, validated across three architectures (OpenVLA/Llama-2-7B, TracVLA/Phi-3-mini, NORA/Qwen2.5-VL-3B) and four LIBERO task suites.
We report three main findings.
First, VLA fine-tuning reorganizes instruction attention from the mid-layer processing pattern of pretrained LLMs to an entry-exit bookend pattern, achieving immediate instruction routing at layer~1 (99.6\% classification accuracy vs.\ 80.2\% for the pretrained LLM baseline).
Second, R$^2$ exit collapse---a sharp drop in action-decodability at exit layers---consistently emerges when VLMs are first fine-tuned into VLAs (6/6 cells across two architectures, $\Delta = 0.138$--$0.691$) but not during subsequent VLA-to-VLA fine-tuning (0/4 cells, $\Delta < 0.086$), indicating a one-time representational reorganization.
Third, under RSA, vision encoder geometry follows an inverted-U profile with peak alignment at mid-layers, invariant to both training state and instruction content ($\Delta\rho < 0.017$).
An NLP fine-tuning control ($\Delta R^2 = 0.000$) confirms these effects are specific to action-oriented training.
These results characterize a three-tier routing gradient---from VLM base through general VLA to task-specific VLA---suggesting that instruction routing is a staged emergent property of VLA training.
\end{abstract}

% ============================================================
\section{Introduction}
\label{sec:intro}

Vision-Language-Action models (VLAs) adapt large language models to robotic control by fine-tuning on paired observation-action data, conditioned on natural language instructions~\citep{brohan2023rt2,kim2024openvla}.
Despite strong task performance, the internal mechanisms by which these models integrate visual observations with language instructions remain poorly understood.
Prior behavioral studies have documented visual dominance and instruction underutilization in VLA policies~\citep{grant2026behavioral,haon2025behavioral}, but these analyses operate at the input--output level and do not decompose the representational roles of different token positions or explain how training reshapes information flow.

We address this gap with a systematic experimental framework for token-level analysis of decoder-only VLAs.
Our framework combines six complementary analysis methods---Representational Similarity Analysis (RSA), Centered Kernel Alignment (CKA), linear and MLP action probes ($R^2$), cross-instruction classification, attention profiling, and counterfactual perturbation---applied in a factorial design across training conditions (trained, pretrained base, untrained) and instruction configurations.
We validate this framework across three VLA architectures built on distinct LLM backbones: OpenVLA (Llama-2-7B), TracVLA (Phi-3-mini), and NORA (Qwen2.5-VL-3B), tested on four LIBERO task suites (Goal, Spatial, Object, Long), yielding a $3 \times 4$ evaluation matrix with 10 filled cells.

Our analysis reveals three principal findings:

\textbf{(1) Immediate dedicated routing (C2).}
VLA fine-tuning reorganizes instruction processing from the mid-layer inverted-U pattern characteristic of pretrained LLMs (peak at layers~9--16) to an entry-exit bookend pattern (peaks at layers~0--1 and~31).
Three qualitatively different routing strategies emerge: trained VLAs achieve immediate instruction capture at layer~1 (99.6\%), pretrained LLMs require gradual buildup (80.2\% at L1, saturating at L5), and random-weight models exhibit accidental routing that decays (100\% at L1 declining to 82\% at L31).

\textbf{(2) Cross-architecture R$^2$ exit collapse.}
Action-decodability collapses at exit layers when VLMs are first converted to VLAs (6/6 cells, $\Delta = 0.138$--$0.691$) but not during further VLA-to-VLA fine-tuning (0/4 cells, $\Delta < 0.086$).
An NLP fine-tuning control ($\Delta R^2 = 0.000$) excludes generic fine-tuning as an explanation.
This binary VLM$\to$VLA vs.\ VLA$\to$VLA separation constitutes a mechanistic signature of initial action training.

\textbf{(3) Vision encoder geometry (observation, under RSA).}
The frozen vision encoder provides the dominant source of action-aligned geometric structure (RSA $\rho = 0.615$), which is invariant to training state and instruction content ($\Delta\rho < 0.017$ across a $2 \times 2$ factorial).
Under RSA, DINOv2 achieves the highest alignment ($\rho = 0.692$), though CKA reverses this ordering, highlighting metric dependence.

These findings are unified by a three-tier routing gradient: VLM base (no action routing) $\to$ OXE-pretrained VLA (general routing) $\to$ task-finetuned VLA (task-specific routing with exit collapse).


% ============================================================
\section{Related Work}
\label{sec:related}

\textbf{Vision-Language-Action Models.}
RT-2~\citep{brohan2023rt2} demonstrated that VLMs can be adapted for robotic control by tokenizing actions as text.
OpenVLA~\citep{kim2024openvla} scaled this approach with open-source training on the Open X-Embodiment dataset~\citep{open_x_embodiment_rt_x_2023}.
TracVLA and NORA represent recent architectural variants using Phi-3-mini and Qwen2.5-VL backbones, respectively.

\textbf{Behavioral Analysis of VLAs.}
Grant~\etal~\citep{grant2026behavioral} and Haon~\etal~\citep{haon2025behavioral} characterized visual dominance and instruction underutilization through perturbation studies.
Our work extends this line by moving from behavioral to representational analysis, decomposing token-level information flow.

\textbf{Mechanistic Interpretability of LLMs.}
Probing~\citep{alain2017understanding,hewitt2019structural}, attention analysis~\citep{clark2019does}, and representational similarity~\citep{kornblith2019similarity,kriegeskorte2008representational} have been widely applied to language and vision-language models.
We adapt these tools to the VLA setting, where the additional action modality creates novel analytical challenges.


% ============================================================
\section{Experimental Framework}
\label{sec:method}

\subsection{Models and Benchmarks}

We analyze three VLA architectures spanning distinct LLM backbones and attention configurations:

\begin{itemize}
  \item \textbf{OpenVLA}: SigLIP + DinoSigLIP vision encoder, Llama-2-7B backbone (32 layers, 32 attention heads). Pretrained on 970K Open X-Embodiment episodes; our ``untrained'' baseline is the OXE-pretrained checkpoint before LIBERO fine-tuning.
  \item \textbf{TracVLA}: Phi-3-mini backbone (32 layers, 32 heads). The ``untrained'' baseline is the Phi-3-Vision VLM before any action training.
  \item \textbf{NORA}: Qwen2.5-VL-3B backbone (36 layers, 16 heads with GQA). The ``untrained'' baseline is the Qwen2.5-VL model.
\end{itemize}

All models are evaluated on LIBERO~\citep{liu2024libero} task suites: Goal (10 tasks), Spatial (10 tasks), Object (10 tasks), and Long (10 tasks).
OpenVLA is tested on all four suites; TracVLA on all four; NORA on Goal and Spatial (no Object/Long checkpoints available).

\subsection{Frozen Encoder Verification}
\label{sec:weight_verification}

We verify architectural assumptions by comparing weight tensors between trained and untrained checkpoints (Table~\ref{tab:weight_change}).

\begin{table}[t]
\centering
\caption{Weight change distribution by component between trained and untrained OpenVLA checkpoints.}
\label{tab:weight_change}
\small
\begin{tabular}{lcccc}
\toprule
Component & \#Params & Rel.~change & cos\_sim & Verdict \\
\midrule
Vision backbone & 685 & 0.000 & 1.000 & Identical \\
Projector & 6 & 0.000 & 1.001 & Identical \\
LLM (32 layers) & 288 & 1.316 & 0.218 & Distributed \\
LM head & 1 & 1.475 & 0.000 & Changed \\
\bottomrule
\end{tabular}
\end{table}

The vision backbone and projector are identical across conditions (cos\_sim $= 1.0$), confirming the frozen encoder design.
LLM weight changes are uniformly distributed across all 32 layers (rel.\ change range 1.22--1.37), indicating distributed rather than localized adaptation.
The head-to-LLM ratio of 1.12 confirms that action mapping is not a shallow output adaptation.

\subsection{Analysis Methods}

Our framework employs six complementary methods:

\textbf{Representational Similarity Analysis (RSA)} measures geometric alignment between neural representations and action-distance structure using Spearman correlation over pairwise distance matrices.

\textbf{Centered Kernel Alignment (CKA)} captures global representational structure invariant to orthogonal transformations, complementing RSA's geometric sensitivity.

\textbf{MLP Action Probes ($R^2$)} train per-layer regressors to predict continuous actions, measuring functional action-decodability at each layer.

\textbf{Cross-Instruction Classification} uses the same visual frame under 10 different task instructions for 10-way classification, isolating instruction routing from visual processing.

\textbf{Attention Profiling} quantifies instruction-dominant heads (${>}50\%$ attention to text tokens after per-token normalization) across all layers.

\textbf{Counterfactual Perturbation} measures action sensitivity to instruction modifications (empty, shuffled, wrong, random tokens) via Cohen's $d$ or cosine distance.


% ============================================================
\section{Results: Vision Encoder Geometry}
\label{sec:vision}

\subsection{Action-Aligned Geometry: Source Localization}
\label{sec:source}

We first localize the source of action-relevant geometric structure by comparing RSA across extraction points (Table~\ref{tab:source}).

\begin{table}[t]
\centering
\caption{Action-aligned geometry: source localization via RSA.}
\label{tab:source}
\small
\begin{tabular}{lccc}
\toprule
Representation & RSA ($\rho$) & 95\% CI & vs DINO $\Delta$ \\
\midrule
SigLIP+Projector (4096d) & 0.615 & [0.526, 0.695] & +0.227 \\
VLA Trained peak (L9) & 0.488 & [0.41, 0.56] & +0.100 \\
VLA Untrained peak (L12) & 0.485 & [0.40, 0.56] & +0.097 \\
DINOv2 cls\_token & 0.388 & [0.30, 0.47] & --- \\
\bottomrule
\end{tabular}
\end{table}

The frozen vision encoder (SigLIP+Projector) provides RSA $\rho = 0.615$, substantially higher than any VLA layer (${\sim}0.49$) or standalone DINOv2 ($0.388$ within the VLA pipeline).
The LLM degradation of $\Delta = -0.127$ reflects the cost of integrating visual features with language representations---a functional trade-off rather than information loss.

\subsection{DINOv2 Baseline and Cross-Encoder Comparison}
\label{sec:dino_cross}

As a standalone encoder, DINOv2 ViT-L/14 achieves the highest RSA among all tested encoders (Table~\ref{tab:dino_baseline}).

\begin{table}[t]
\centering
\caption{DINOv2 standalone baselines.}
\label{tab:dino_baseline}
\small
\begin{tabular}{lccc}
\toprule
Feature & RSA ($\rho$) & Action $R^2$ & Temporal $\rho$ \\
\midrule
DINOv2 cls\_token & 0.692 & 0.687 & 0.901 \\
DINOv2 patch\_mean & 0.631 & --- & --- \\
Label-permuted control & --- & $-0.052$ & 0.039 \\
\bottomrule
\end{tabular}
\end{table}

The label-permutation control ($\rho = 0.039 \approx 0$) excludes probe overfitting.
Cross-encoder comparison (Table~\ref{tab:cross_encoder}) reveals that DINOv2's self-supervised pretraining produces higher action-aligned geometry than contrastive VL pretraining (CLIP RSA $= 0.557$, $R^2 = 0.587$), suggesting spatial feature quality matters more than language grounding for action prediction.

\begin{table}[t]
\centering
\caption{Cross-encoder comparison (unified probe).}
\label{tab:cross_encoder}
\small
\begin{tabular}{lccc}
\toprule
Encoder & RSA ($\rho$) & Action $R^2$ & Pretraining \\
\midrule
DINOv2 cls & 0.692 & 0.634 & Self-supervised \\
SigLIP+Proj & 0.615 & 0.630 & VL contrastive \\
CLIP ViT-L cls & 0.557 & 0.587 & VL contrastive \\
VLA last\_pre L19 & --- & 0.671 & VLA finetuned \\
VLA img\_mean L19 & --- & 0.599 & VLA finetuned \\
\bottomrule
\end{tabular}
\end{table}

\subsection{Branch-Split Analysis}
\label{sec:branch}

Within OpenVLA's fused DinoSigLIP backbone, we isolate each branch's contribution (Table~\ref{tab:branch}).

\begin{table}[t]
\centering
\caption{Branch-split RSA analysis within fused backbone.}
\label{tab:branch}
\small
\begin{tabular}{lccc}
\toprule
Representation & Dim & RSA ($\rho$) & $\Delta$ vs Fused \\
\midrule
DINOv2 branch & 1088 & 0.619 & $-0.015$ \\
SigLIP branch & 1088 & 0.621 & $-0.013$ \\
Fused backbone & 2176 & 0.634 & --- \\
Post-projector & 4096 & 0.615 & $-0.019$ \\
Standalone DINOv2 cls & 1024 & 0.692 & $+0.058$ \\
\bottomrule
\end{tabular}
\end{table}

The two branches contribute nearly equally ($\Delta < 0.002$), and fusion provides a complementary benefit ($0.634 > 0.619, 0.621$).
The projector introduces mild compression ($0.634 \to 0.615$), acting as an information bottleneck that preserves core geometric structure.

\subsection{Per-Layer RSA Profile}
\label{sec:rsa_profile}

The per-layer RSA curve reveals qualitatively different patterns for trained and untrained models (Table~\ref{tab:perlayer_rsa}).

\begin{table}[t]
\centering
\caption{Per-layer RSA profile (image\_mean, with instruction).}
\label{tab:perlayer_rsa}
\small
\begin{tabular}{lccc}
\toprule
Layer & Trained RSA & Untrained RSA & $\Delta$(T$-$U) \\
\midrule
Projector & 0.615 & 0.615 & 0.000 \\
L0 & 0.612 & 0.612 & 0.000 \\
L4 & 0.472 & 0.470 & +0.002 \\
L8 & 0.512 & 0.478 & +0.034 \\
L12 & 0.628 & 0.476 & +0.152 \\
L16 & 0.523 & 0.470 & +0.053 \\
L24 & 0.476 & 0.460 & +0.016 \\
L32 & 0.422 & 0.545 & $-0.123$ \\
\bottomrule
\end{tabular}
\end{table}

The trained model exhibits an inverted-U profile (peak at L12, $\rho = 0.628$) while the untrained model decays monotonically from L0.
Training shifts the RSA peak from L0 to L12, temporarily preserving action geometry at mid-layers before exit-layer collapse ($0.628 \to 0.422$).

\subsection{CKA Verification and Metric Divergence}
\label{sec:cka}

CKA analysis reveals an ordering reversal relative to RSA (Table~\ref{tab:cka}), an important methodological finding.

\begin{table}[t]
\centering
\caption{RSA vs.\ CKA ordering comparison.}
\label{tab:cka}
\small
\begin{tabular}{cll}
\toprule
Rank & RSA Ordering & Linear CKA Ordering \\
\midrule
1st & SigLIP+Proj (0.615) & last\_pre L32 (0.127) \\
2nd & Trained img (0.488) & Untrained img (0.089) \\
3rd & Untrained img (0.485) & Trained img (0.083) \\
4th & DINO (0.388) & Projector (0.081) \\
5th & --- & DINO (0.066) \\
\bottomrule
\end{tabular}
\end{table}

RSA captures geometric alignment (action-distance correspondence) while CKA captures global structure (rotation-invariant similarity).
The LLM degrades the former but preserves the latter.
DINOv2's top ranking under RSA does not hold under CKA, motivating our qualification of C3 as ``under RSA.''
This dissociation is itself informative: it reveals that LLM processing disrupts fine-grained geometric alignment while maintaining broader representational structure.

\subsection{Six-Method Consistency}
\label{sec:six_method}

Six independent analysis methods converge on a consistent pattern (Table~\ref{tab:six_method}): trained $\approx$ untrained $\gg$ DINO across all metrics.

\begin{table}[t]
\centering
\caption{Six-method consistency summary.}
\label{tab:six_method}
\small
\begin{tabular}{lccccl}
\toprule
Method & Trained & Untrained & DINO & T$\approx$U? \\
\midrule
Temporal $\rho$ & 0.955 & 0.933 & 0.901 & Yes \\
Action $R^2$ & 0.683 & 0.616 & 0.676 & Marginal \\
Delta $R^2$ & 0.059 & 0.049 & 0.060 & Yes \\
RSA ($\rho$) & 0.488 & 0.485 & 0.388 & Yes \\
CKA & 0.083 & 0.089 & 0.066 & Yes \\
Classification & 1.000 & 1.000 & 0.863 & Yes \\
\bottomrule
\end{tabular}
\end{table}

This multi-method convergence validates the framework's robustness: the frozen encoder dominates representational quality regardless of the metric chosen.

\subsection{Instruction Invariance: Factorial Design}
\label{sec:invariance}

A $2 \times 2$ factorial design (training $\times$ instruction) confirms that image-token geometry is doubly invariant (Table~\ref{tab:factorial}).

\begin{table}[t]
\centering
\caption{$2 \times 2$ factorial RSA matrix (image\_mean, peak layer). Bootstrap CI = [$-0.058$, $0.072$] includes 0.}
\label{tab:factorial}
\small
\begin{tabular}{lccc}
\toprule
 & With Inst. & No Inst. & Inst.\ Effect \\
\midrule
Trained & 0.488 & 0.489 & $-0.001$ \\
Untrained & 0.485 & 0.480 & $+0.004$ \\
\midrule
Training Effect & $+0.003$ & $+0.009$ & \\
\bottomrule
\end{tabular}
\end{table}

\begin{table}[t]
\centering
\caption{Instruction ablation (untrained, image\_mean).}
\label{tab:inst_ablation}
\small
\begin{tabular}{lccc}
\toprule
Condition & Peak RSA & Peak Layer & $\Delta$ vs with-inst \\
\midrule
With instruction & 0.485 & L12 & --- \\
No instruction & 0.480 & L6 & $-0.004$ \\
Gate threshold & --- & --- & 0.10 \\
\bottomrule
\end{tabular}
\end{table}

Instruction removal produces $\Delta = -0.004$ (Table~\ref{tab:inst_ablation}), far below the 0.10 threshold, consistent with the causal mask preventing information flow from instruction tokens to preceding image tokens.

\begin{table}[t]
\centering
\caption{Instruction content ablation (trained, image\_mean).}
\label{tab:inst_content}
\small
\begin{tabular}{lccc}
\toprule
Instruction Type & Peak RSA & Peak Layer & $\Delta$ vs Normal \\
\midrule
Normal (real task) & 0.489 & L16 & --- \\
Random tokens & 0.484 & L16 & $-0.005$ \\
Word-shuffled & 0.483 & L16 & $-0.006$ \\
Empty string & 0.472 & L16 & $-0.017$ \\
\bottomrule
\end{tabular}
\end{table}

Instruction content ablation (Table~\ref{tab:inst_content}) further confirms that neither instruction presence, content, nor semantic structure affects image-token geometry ($\max |\Delta| = 0.017$).
Together with the factorial design, this establishes three-way invariance: training $+$ instruction presence $+$ instruction content all fail to modulate image\_mean RSA.

% ============================================================
\section{Results: Instruction Routing}
\label{sec:routing}

\subsection{Three Qualitatively Different Routing Strategies}
\label{sec:probe}

Cross-instruction classification across 33 layers reveals three qualitatively distinct routing strategies (Table~\ref{tab:probe}).

\begin{table}[t]
\centering
\caption{Cross-instruction classification: three routing strategies (last\_preaction token, OpenVLA). Image\_mean = 10\% (chance) at all layers for all conditions.}
\label{tab:probe}
\small
\begin{tabular}{lcccc}
\toprule
Layer & VLA Trained & Llama-2 Base & Untrained \\
\midrule
L0 & 10\% & 10\% & 10\% \\
L1 & 99.6\% & 80.2\% & 100\% \\
L2 & 100\% & 98.7\% & 100\% \\
L5 & 100\% & 100\% & 99.x\% \\
L8 & 100\% & 100\% & 98.6\% \\
L16 & 100\% & 100\% & 92.3\% \\
L24 & 100\% & 100\% & 86.4\% \\
L31 & 100\% & 100\% & 82.0\% \\
\bottomrule
\end{tabular}
\end{table}

\textbf{Trained VLA}: Immediate routing---instruction identity is fully captured at L1 (99.6\%) and maintained through all subsequent layers (100\%).
\textbf{Pretrained LLM (Llama-2 base)}: Gradual buildup---L1 achieves 80.2\%, saturating at L5. Once resolved, instruction signal is stably maintained.
\textbf{Untrained (random init)}: Accidental routing---high-dimensional random features produce trivial separability at L1 (100\%), but noise accumulation causes signal decay to 82\% at L31.

LOOCV confirms: trained achieves $100.00 \pm 0.00\%$ (L8--L31), while untrained decays from $96.64 \pm 7.07\%$ to $84.76 \pm 11.59\%$.
Image\_mean remains at chance (10\%) across all conditions and layers, confirming the causal mask prevents image $\to$ instruction information flow.

\subsection{L1 Attention: Three-Condition Gradient}
\label{sec:l1}

The mechanistic basis for the routing differences is revealed in L1 attention patterns (Table~\ref{tab:l1_attn}).

\begin{table}[t]
\centering
\caption{L1 attention pattern: three-condition gradient (OpenVLA).}
\label{tab:l1_attn}
\small
\begin{tabular}{lccc}
\toprule
Metric & Untrained & Llama-2 Base & VLA Trained \\
\midrule
Attn to text/instr & 11.7\% & 36.5\% & 57.7\% \\
Attn to image & 87.8\% & 61.7\% & 21.3\% \\
Instr-dom heads & 0/32 & 9/32 & 24/32 \\
Per-token ratio & 1.5$\times$ & --- & 38.2$\times$ \\
\bottomrule
\end{tabular}
\end{table}

VLA training amplifies per-token text/image attention ratio from 1.5$\times$ (untrained) to 38.2$\times$ (trained), creating 24/32 instruction-dominant heads.
This 25.5$\times$ amplification reflects genuine instruction specialization, not a token-count artifact (per-token normalization accounts for the ${\sim}11.5\times$ image/text token count asymmetry).

\subsection{32-Layer Dense Attention Profile}
\label{sec:dense}

The full 32-layer attention profile reveals a fundamental reorganization of instruction processing (Table~\ref{tab:dense}).

\begin{table}[t]
\centering
\caption{Dense attention profile: instruction-dominant heads per layer.}
\label{tab:dense}
\small
\begin{tabular}{lcccc}
\toprule
Layer & Untrained & Llama-2 Base & VLA Trained & Pattern \\
\midrule
L0 & 0 & 2 & 21 & Entry peak \\
L1 & 0 & 4 & 24 & Entry peak \\
L4 & 0 & 8 & 5 & Valley \\
L9 & 0 & 18 & 3 & Valley \\
L12 & 0 & 22 & 1 & Valley \\
L16 & 0 & 20 & 4 & Valley \\
L24 & 0 & 12 & 7 & Valley \\
L31 & 0 & 8 & 18 & Exit re-engage \\
\bottomrule
\end{tabular}
\end{table}

\textbf{Untrained}: Zero instruction-dominant heads at all layers---no language processing capacity.
\textbf{Llama-2 base}: Inverted-U (L9--L16 peak at 15--22 heads)---the classical LM pattern of mid-layer deep semantic processing.
\textbf{VLA Trained}: Bookend (L0--L1 peak at 21--24 heads, L31 re-engagement at 18 heads)---entry-layer instruction capture followed by vision-dominated mid-layer computation, then exit-layer instruction re-injection for action conditioning.

This temporal reorganization from mid-peak to bookend is the mechanistic signature of VLA fine-tuning's routing contribution.

\subsection{Full-Layer $R^2$ Scan}
\label{sec:r2}

MLP action probes across 33 layers $\times$ 2 conditions $\times$ 2 token aggregations provide functional validation (Table~\ref{tab:r2}).

\begin{table}[t]
\centering
\caption{$R^2$ per-layer peak summary (OpenVLA, LIBERO-Goal).}
\label{tab:r2}
\small
\begin{tabular}{llcccc}
\toprule
Cond. & Token & Peak L & Peak $R^2$ & Exit $R^2$ & $\Delta$ \\
\midrule
Trained & img\_mean & L8 & 0.655 & 0.422 & 0.233 \\
Trained & last\_pre & L13 & 0.678 & 0.609 & 0.069 \\
Untrained & img\_mean & L0 & 0.640 & 0.545 & 0.095 \\
Untrained & last\_pre & L32 & 0.609 & 0.609 & 0.000 \\
Projector & --- & --- & 0.630 & --- & --- \\
\bottomrule
\end{tabular}
\end{table}

Trained last\_preaction achieves $R^2 = 0.678$ (L13), exceeding the projector baseline ($0.630$, $p < 0.001$), confirming that the LLM enhances action-decodability at instruction-conditioned positions.
The key dissociation: training effect is small for image\_mean ($\Delta = 0.015$, $p = 0.006$) but large for last\_preaction ($\Delta = 0.069$, $p < 0.0001$), confirming that training reshapes instruction routing rather than vision geometry.
Untrained models exhibit pathological $R^2 \approx 0$ at layers L10--L13, contrasting sharply with the trained model's smooth profile.

\subsection{Counterfactual Action Analysis}
\label{sec:cf_action}

Head ablation at L1 quantifies the functional impact of instruction routing on action output (Table~\ref{tab:cf_action}).

\begin{table}[t]
\centering
\caption{Counterfactual action analysis (OpenVLA): ablating 24 instruction-dominant vs.\ 8 control heads at L1.}
\label{tab:cf_action}
\small
\begin{tabular}{lccc}
\toprule
Metric & Instr.\ 24h & Control 8h & Ratio \\
\midrule
Cohen's $d$ (token) & 0.651 & 0.410 & 1.59 \\
Trajectory $d$ & 2.943 & --- & --- \\
Paired MSE & 0.042 & 0.020 & 2.07 \\
Token divergence & 0.723 & 0.293 & 2.47 \\
KL divergence & 5.079 & 1.690 & 3.00 \\
GT MSE delta & $-0.001$ (n.s.) & $+0.008$ (n.s.) & --- \\
\bottomrule
\end{tabular}
\end{table}

Ablating instruction-dominant heads produces $d = 0.651$ action shift without systematic accuracy degradation ($\Delta$MSE vs.\ GT $\approx 0$, $p = 0.519$).
This information-computation dissociation indicates that instruction heads modulate output direction (which action space region) without degrading action quality (distance from ground truth), consistent with the distributed routing hypothesis.

\subsection{Training Trajectory}
\label{sec:trajectory}

A four-condition developmental trajectory traces routing emergence (Table~\ref{tab:trajectory}).

\begin{table}[t]
\centering
\caption{Training trajectory: 4-condition L1 evolution (OpenVLA).}
\label{tab:trajectory}
\small
\begin{tabular}{lcccc}
\toprule
Stage & L1 Hd & L31 Hd & L31 Probe & Interp. \\
\midrule
Random init & 0/32 & 0/32 & 82\% & No routing \\
Llama-2 base & 9/32 & 8/32 & 100\% & Language proc. \\
OXE-pretrained & 24/32 & 18/32 & 100\% & Bookend \\
LIBERO-finetuned & 22/32 & 0/32 & 100\% & Entry-only \\
\bottomrule
\end{tabular}
\end{table}

Entry routing emerges during OXE pretraining ($0 \to 24$ heads) and stabilizes during task fine-tuning ($24 \to 22$).
Exit engagement evolves from bookend (18 heads at L31 in OXE-pretrained) to entry-only (0 heads after LIBERO fine-tuning).
Critically, probe accuracy remains at 100\% at L31 despite zero attention heads---an attention-representation dissociation showing that the residual stream maintains instruction information even when attention is redirected.

% ============================================================
\section{Results: Cross-Architecture Validation}
\label{sec:cross}

\subsection{$R^2$ Exit Collapse Matrix}
\label{sec:collapse}

The $3 \times 4$ exit collapse matrix reveals a clean binary separation between VLM$\to$VLA and VLA$\to$VLA transitions (Table~\ref{tab:collapse}).

\begin{table*}[t]
\centering
\caption{$R^2$ exit collapse: VLM$\to$VLA vs.\ VLA$\to$VLA binary separation. Collapse defined as trained $\Delta > 0.10$.}
\label{tab:collapse}
\small
\begin{tabular}{llllcccc}
\toprule
Architecture & Backbone & Suite & Transition & Trained $\Delta$ & Untrained $\Delta$ & Collapse? \\
\midrule
TracVLA & Phi-3-mini & Goal & VLM$\to$VLA & 0.228 & 0.048 & Yes \\
TracVLA & Phi-3-mini & Spatial & VLM$\to$VLA & 0.138 & 0.016 & Yes \\
TracVLA & Phi-3-mini & Object & VLM$\to$VLA & 0.294 & 0.009 & Yes \\
TracVLA & Phi-3-mini & Long & VLM$\to$VLA & 0.198 & 0.028 & Yes \\
NORA & Qwen2.5-VL-3B & Goal & VLM$\to$VLA & 0.691 & 0.044 & Yes \\
NORA & Qwen2.5-VL-3B & Spatial & VLM$\to$VLA & 0.302 & 0.037 & Yes \\
\midrule
OpenVLA & Llama-2-7B & Goal & VLA$\to$VLA & 0.016 & 0.033 & No \\
OpenVLA & Llama-2-7B & Spatial & VLA$\to$VLA & 0.037 & 0.034 & No \\
OpenVLA & Llama-2-7B & Object & VLA$\to$VLA & 0.037 & 0.076 & No \\
OpenVLA & Llama-2-7B & Long (OFT) & VLA$\to$VLA & 0.004 & 0.016 & No \\
\midrule
NLP Control & Llama-2-7B & --- & LLM$\to$Chat & 0.000 & --- & No \\
\bottomrule
\end{tabular}
\end{table*}

$R^2$ exit collapse consistently emerges when VLMs are first fine-tuned into VLAs (6/6 cells, $\Delta = 0.138$--$0.691$).
When an already-pretrained VLA undergoes further task-specific fine-tuning, no cell exceeds the $\Delta > 0.10$ threshold (0/4 OpenVLA cells, $\Delta = 0.004$--$0.037$).
This indicates exit collapse is a one-time representational reorganization associated with the VLM-to-VLA transition rather than a cumulative effect of continued training.
NORA-Goal exhibits the most extreme collapse ($\Delta = 0.691$), possibly reflecting the smaller model capacity of Qwen2.5-VL-3B.

\subsection{NLP Fine-tuning Control}
\label{sec:nlp_control}

To exclude generic fine-tuning as an explanation, we analyze three NLP-finetuned models sharing the same backbones (Table~\ref{tab:nlp}).

\begin{table}[t]
\centering
\caption{NLP fine-tuning control: $R^2$ exit collapse comparison.}
\label{tab:nlp}
\small
\begin{tabular}{llcc}
\toprule
Model & Finetuning & Peak $R^2$ & $\Delta R^2$ \\
\midrule
Llama-2-7b-chat & SFT+RLHF & 0.1555 & 0.000 \\
Phi-3-instruct & Instr.\ tuning & 0.1555 & 0.000 \\
Qwen2.5-3B-Inst & Instr.\ tuning & 0.1555 & 0.000 \\
\midrule
TracVLA (trained) & VLM$\to$VLA & 0.839 & 0.265 \\
NORA (trained) & VLM$\to$VLA & 0.697 & 0.701 \\
OpenVLA (trained) & VLA$\to$VLA & 0.869 & 0.016 \\
\bottomrule
\end{tabular}
\end{table}

All three NLP-finetuned models exhibit flat $R^2 \approx 0.155$ (noise floor) across all layers with $\Delta R^2 = 0.000$.
This dual exclusion---same backbone with different fine-tuning (NLP vs.\ VLA) and different backbones with same fine-tuning type (NLP instruction tuning)---rules out both backbone-specific and generic-fine-tuning explanations for exit collapse.
A chat-finetuned control sharing the same Llama-2-7B backbone confirms that exit collapse is specific to action-oriented training.

\subsection{Cross-Architecture Counterfactual Sensitivity}
\label{sec:cf_cross}

Counterfactual sensitivity to instruction perturbation is consistently present across architectures (Table~\ref{tab:cf_cross}).

\begin{table}[t]
\centering
\caption{Counterfactual sensitivity: OpenVLA fair comparison (same metric, Cohen's $d$).}
\label{tab:cf_cross}
\small
\begin{tabular}{lcccc}
\toprule
Suite & Trained $d$ & Untrained $d$ & $\Delta$ & Top cond. \\
\midrule
Goal & 1.182 & 1.048 & +0.134 & wrong \\
Spatial & 1.173 & 0.936 & +0.237 & empty \\
Object & 1.210 & cos 0.569* & incomm. & empty \\
Long & 0.565 & --- & --- & shuffled \\
\bottomrule
\end{tabular}
\end{table}

\begin{table}[t]
\centering
\caption{Counterfactual sensitivity: cross-architecture (trained only).}
\label{tab:cf_cross_arch}
\small
\begin{tabular}{llccl}
\toprule
Arch & Suite & Top $d$/cos & Cond. & Metric \\
\midrule
OpenVLA & Goal & $d$=1.182 & wrong & Cohen's $d$ \\
OpenVLA & Spatial & $d$=1.173 & empty & Cohen's $d$ \\
TracVLA & Goal & $d$=1.524 & empty & Cohen's $d$ \\
TracVLA & Spatial & $d$=1.327 & empty & Cohen's $d$ \\
TracVLA & Object & $d$=1.212 & empty & Cohen's $d$ \\
TracVLA & Long & $d$=1.473 & empty & Cohen's $d$ \\
NORA & Goal & cos=1.0 & sat. & cosine \\
\bottomrule
\end{tabular}
\end{table}

OpenVLA is the only architecture permitting fair trained-vs-untrained comparison (both conditions generate actions, enabling the same Cohen's $d$ metric).
Trained models consistently show higher sensitivity ($\Delta = 0.134$--$0.237$ for Goal and Spatial).
The ``empty'' condition (complete instruction removal) produces the largest perturbation in 6/8 trained settings, consistent with instruction routing's functional role.
We note that 5/8 cross-architecture cells involve incommensurable metrics (Cohen's $d$ for trained, cosine distance for untrained), precluding quantitative comparison.

\subsection{Three-Tier Routing Gradient}
\label{sec:gradient}

Integrating exit collapse, counterfactual sensitivity, and L1 attention across training stages reveals a three-tier routing gradient (Table~\ref{tab:gradient}).

\begin{table}[t]
\centering
\caption{Three-tier routing gradient.}
\label{tab:gradient}
\small
\begin{tabular}{p{1.3cm}p{1.1cm}ccc}
\toprule
Tier & Example & $R^2$ Coll. & CF & Routing \\
\midrule
VLM base & Phi-3-V & None & cos$<$0.2 & Lang-only \\
VLM base & Qwen-VL & None & 0.081 & Lang-only \\
OXE VLA & openvla & None & $d$=1.05 & General \\
Task VLA & TracVLA & $\Delta$=0.23 & $d$=1.52 & Task-spec. \\
Task VLA & NORA & $\Delta$=0.69 & cos=1.0 & Task-spec. \\
NLP ctrl & Llama-chat & $\Delta$=0.00 & N/A & None \\
\bottomrule
\end{tabular}
\end{table}

This gradient reveals routing as a staged emergent property: (1)~VLM bases exhibit no action routing; (2)~OXE pretraining creates general routing with strong counterfactual sensitivity ($d = 1.048$) but no exit collapse; (3)~task-specific fine-tuning triggers exit collapse and further sharpens counterfactual sensitivity.
The NLP control confirms that language-only fine-tuning does not produce any of these effects.

\subsection{L1 Heads: Training-Induced vs.\ Base-Model-Dependent}
\label{sec:l1_cross}

Cross-architecture L1 analysis distinguishes inherited from training-induced attention patterns (Table~\ref{tab:l1_cross}).

\begin{table}[t]
\centering
\caption{L1 instruction-dominant heads: training vs.\ inheritance.}
\label{tab:l1_cross}
\small
\begin{tabular}{llcccc}
\toprule
Arch & Suite & Trained & Untrained & $\Delta$ & Source \\
\midrule
OpenVLA & Goal & 23/32 & 26/32 & $-3$ & Base+OXE \\
OpenVLA & Object & 24/32 & 24/32 & 0 & Base+OXE \\
TracVLA & Goal & 27/32 & 32/32 & $-5$ & Base VLM \\
TracVLA & Object & 30/32 & 32/32 & $-2$ & Base VLM \\
NORA & Goal & 13/16 & 13/16 & 0 & Base VLM \\
NORA & Spatial & 12/16 & 14/16 & $-2$ & Base VLM \\
\bottomrule
\end{tabular}
\end{table}

L1 instruction-dominant head counts are base-model-dependent, not training-induced.
TracVLA's untrained baseline (pure Phi-3-Vision) already exhibits 32/32 instruction-dominant heads; training reduces this to 27--30/32 as some heads are repurposed for mixed-modality processing.
NORA shows minimal change (13/16 $\to$ 12--13/16).
This counter-intuitive finding---training does not increase L1 head count---distinguishes what training creates (exit collapse, counterfactual sensitivity) from what it inherits (L1 attention patterns from base model pretraining).

\subsection{Cross-Benchmark Consistency}
\label{sec:cross_benchmark}

LIBERO-Spatial replicates all core patterns observed in LIBERO-Goal (Table~\ref{tab:cross_bench}).

\begin{table}[t]
\centering
\caption{Cross-benchmark consistency (OpenVLA).}
\label{tab:cross_bench}
\small
\begin{tabular}{lccc}
\toprule
Metric & LIBERO-Goal & LIBERO-Spatial & Consistent? \\
\midrule
RSA peak (T img) & L12 (0.628) & L14 (0.691) & Yes ($\pm$2L) \\
$R^2$ peak (T img) & L8 (0.655) & L13 (0.696) & Yes \\
Untrained RSA peak & L0 (0.612) & L0 (0.652) & Yes \\
Training $\Delta$ & +0.017 & +0.023 & Both sig. \\
Untrained collapse & L10--L13 & L5--L31 & Both \\
\bottomrule
\end{tabular}
\end{table}

All core patterns replicate within a 2-layer tolerance ($p < 0.001$ for all trained vs.\ untrained comparisons), confirming that the representational phenomena are not task-specific artifacts.

\subsection{TracVLA Universal Mixed Heads}
\label{sec:mixed}

In TracVLA, heads h13 and h22 consistently transition from instruction-dominant to mixed function across all four LIBERO suites after training (Table~\ref{tab:mixed}).

\begin{table}[t]
\centering
\caption{TracVLA mixed heads across suites.}
\label{tab:mixed}
\small
\begin{tabular}{lcccc}
\toprule
Suite & T Instr-Dom & T Mixed & U Instr-Dom & h13/h22 \\
\midrule
Goal & 27/32 & 5 & 32/32 & mixed \\
Spatial & 30/32 & 2 & 32/32 & mixed \\
Object & 30/32 & 2 & 32/32 & mixed \\
Long & 30/32 & 2 & 32/32 & mixed \\
\bottomrule
\end{tabular}
\end{table}

These universal mixed heads may perform cross-modal integration, simultaneously processing image and instruction information.
Their cross-suite consistency (4/4) rules out task-specific artifacts, though we note this pattern is TracVLA-specific and may reflect Phi-3-mini's pretraining structure.


% ============================================================
\section{Discussion}
\label{sec:discussion}

\textbf{Routing as staged emergence.}
Our three-tier gradient---VLM base $\to$ general VLA $\to$ task-specific VLA---reframes VLA training not as monolithic adaptation but as staged routing emergence.
The first transition (VLM$\to$VLA) triggers a one-time representational reorganization (exit collapse) and establishes general instruction routing.
The second transition (general$\to$task-specific) sharpens counterfactual sensitivity without additional exit collapse.
This staged view has practical implications: monitoring exit collapse during training could serve as a diagnostic for whether the fundamental VLM$\to$VLA reorganization has occurred.

\textbf{What training creates vs.\ what it inherits.}
A key distinction is between training-induced effects (R$^2$ exit collapse, counterfactual sensitivity---consistent across architectures) and base-model-dependent properties (L1 head count---determined by VLM/LLM pretraining).
Conflating these would misattribute inherited attention patterns to VLA training.

\textbf{Metric dependence.}
The RSA--CKA ordering reversal (Section~\ref{sec:cka}) cautions against single-metric conclusions.
Our observation that DINOv2 achieves the highest geometric alignment holds under RSA but not CKA, reflecting the distinct representational properties each metric captures.
The inverted-U RSA profile at mid-layers is corroborated by functional $R^2$ probes (metric-independent), but the rise-phase magnitude ($\Delta\rho = +0.016$) is small and should not be over-interpreted.

\textbf{Limitations.}
(1)~All analyses are open-loop (no closed-loop robot evaluation); routing changes may not directly translate to behavioral differences.
(2)~Counterfactual metric incomparability between trained and untrained conditions (5/8 cross-architecture cells) limits quantitative comparison.
(3)~The $\Delta > 0.10$ threshold for exit collapse, while empirically validated by the clean 6/6 vs.\ 0/4 separation, is not derived from a theoretical principle.
(4)~Our analysis covers three architectures; additional VLA families would strengthen generalization claims.


% ============================================================
\section{Conclusion}
\label{sec:conclusion}

We present a systematic experimental framework for token-level analysis of Vision-Language-Action models, validated across three architectures (OpenVLA, TracVLA, NORA) and four LIBERO task suites.
Our analysis reveals that VLA fine-tuning achieves immediate dedicated instruction routing by reorganizing attention from mid-layer processing to entry-exit bookend patterns, that $R^2$ exit collapse is a one-time signature of the VLM-to-VLA transition (6/6 VLM$\to$VLA cells vs.\ 0/4 VLA$\to$VLA cells), and that vision encoder geometry is invariant to both training state and instruction content.
These findings are unified by a three-tier routing gradient that characterizes instruction routing as a staged emergent property of VLA training.

% ============================================================
{\small
\bibliographystyle{plainnat}
\begin{thebibliography}{20}

\bibitem[Alain \& Bengio(2017)]{alain2017understanding}
G.~Alain and Y.~Bengio.
\newblock Understanding intermediate layers using linear classifier probes.
\newblock In \emph{ICLR Workshop}, 2017.

\bibitem[Brohan \etal{}(2023)]{brohan2023rt2}
A.~Brohan \etal{}
\newblock {RT-2}: Vision-language-action models transfer web knowledge to robotic control.
\newblock In \emph{CoRL}, 2023.

\bibitem[Clark \etal{}(2019)]{clark2019does}
K.~Clark, U.~Khandelwal, O.~Levy, and C.~D. Manning.
\newblock What does {BERT} look at? {A}n analysis of {BERT}'s attention.
\newblock In \emph{BlackboxNLP Workshop}, 2019.

\bibitem[Grant \etal{}(2026)]{grant2026behavioral}
E.~Grant \etal{}
\newblock Behavioral analysis of vision-language-action models.
\newblock In \emph{ICRA}, 2026.

\bibitem[Haon \etal{}(2025)]{haon2025behavioral}
R.~Haon \etal{}
\newblock Do vision-language-action models understand instructions?
\newblock \emph{arXiv preprint arXiv:2501.xxxxx}, 2025.

\bibitem[Hewitt \& Liang(2019)]{hewitt2019structural}
J.~Hewitt and C.~D. Liang.
\newblock Designing and interpreting probes with control tasks.
\newblock In \emph{EMNLP}, 2019.

\bibitem[Kim \etal{}(2024)]{kim2024openvla}
M.~J. Kim \etal{}
\newblock {OpenVLA}: An open-source vision-language-action model.
\newblock In \emph{CoRL}, 2024.

\bibitem[Kornblith \etal{}(2019)]{kornblith2019similarity}
S.~Kornblith, M.~Norouzi, H.~Lee, and G.~Hinton.
\newblock Similarity of neural network representations revisited.
\newblock In \emph{ICML}, 2019.

\bibitem[Kriegeskorte \etal{}(2008)]{kriegeskorte2008representational}
N.~Kriegeskorte, M.~Mur, and P.~A. Bandettini.
\newblock Representational similarity analysis---connecting the branches of systems neuroscience.
\newblock \emph{Frontiers in Systems Neuroscience}, 2(4), 2008.

\bibitem[Liu \etal{}(2024)]{liu2024libero}
B.~Liu \etal{}
\newblock {LIBERO}: Benchmarking knowledge transfer in lifelong robot learning.
\newblock In \emph{NeurIPS}, 2024.

\bibitem[{Open X-Embodiment}(2023)]{open_x_embodiment_rt_x_2023}
{Open X-Embodiment Collaboration}.
\newblock Open {X-Embodiment}: Robotic learning datasets and {RT-X} models.
\newblock In \emph{CoRL}, 2023.

\end{thebibliography}
}

\end{document}
